% =============================================================================
% MVA Graphs in Machine Learning - Common Beamer Configuration
% Michal Valko (Inria & ENS Paris-Saclay)
% 2025 Modern Edition
% =============================================================================

\usetheme{default}
\useoutertheme{infolines}

% -----------------------------------------------------------------------------
% Core packages
% -----------------------------------------------------------------------------
\usepackage{array}
\usepackage{iftex}
\usepackage{xcolor}
\usepackage{graphicx}
\usepackage{standalone}
\usepackage{amsmath,amssymb}
\usepackage{dsfont}
\usepackage{booktabs}
\usepackage{hyperref}

% -----------------------------------------------------------------------------
% Modern Font Configuration
% Uses TeX Gyre Heros (Helvetica-like) which is bundled with TeX Live
% Falls back gracefully if fonts aren't available
% -----------------------------------------------------------------------------
\ifluatex
  \usepackage{fontspec}
  \IfFontExistsTF{Fira Sans}{%
    \setsansfont{Fira Sans}[
      UprightFont = *-Regular,
      BoldFont = *-SemiBold,
      ItalicFont = *-Italic,
      BoldItalicFont = *-SemiBoldItalic,
      Scale = 0.95
    ]
  }{%
    % Fallback to TeX Gyre Heros (bundled with TeX Live)
    \setsansfont{TeX Gyre Heros}[Scale = 0.95]
  }
  \IfFontExistsTF{Fira Code}{%
    \setmonofont{Fira Code}[Scale = 0.85]
  }{%
    % Fallback to TeX Gyre Cursor (bundled with TeX Live)
    \setmonofont{TeX Gyre Cursor}[Scale = 0.85]
  }
\else\ifxetex
  \usepackage{fontspec}
  \IfFontExistsTF{Fira Sans}{%
    \setsansfont{Fira Sans}[Scale = 0.95]
  }{%
    % Fallback to TeX Gyre Heros (bundled with TeX Live)
    \setsansfont{TeX Gyre Heros}[Scale = 0.95]
  }
  \IfFontExistsTF{Fira Code}{%
    \setmonofont{Fira Code}[Scale = 0.85]
  }{%
    % Fallback to TeX Gyre Cursor (bundled with TeX Live)
    \setmonofont{TeX Gyre Cursor}[Scale = 0.85]
  }
\fi\fi

% -----------------------------------------------------------------------------
% TikZ and graphics
% -----------------------------------------------------------------------------
\usepackage{tikz}
\usepackage{etex}
\usepackage{circuitikz}

% Graph drawing requires LuaLaTeX
\ifluatex
  \usetikzlibrary{graphs,graphdrawing,quotes,shapes,arrows,positioning,graphs.standard,mindmap,calc}
  \usegdlibrary{layered,trees,force}
\else
  \usetikzlibrary{graphs,quotes,shapes,arrows,positioning,graphs.standard,mindmap,calc}
\fi

\usetikzlibrary{overlay-beamer-styles}

% TikZ slide animation helpers
\tikzset{onslide/.code args={<#1>#2}{%
  \only<#1>{\pgfkeysalso{#2}}
}}
\tikzset{temporal/.code args={<#1>#2#3#4}{%
  \temporal<#1>{\pgfkeysalso{#2}}{\pgfkeysalso{#3}}{\pgfkeysalso{#4}}
}}
\tikzstyle{highlight}=[red,ultra thick]

% -----------------------------------------------------------------------------
% Bibliography
% -----------------------------------------------------------------------------
\usepackage[style=apa,natbib=true,backend=biber]{biblatex}
\usepackage{silence}
\WarningFilter{biblatex}{Patching footnotes failed}
\bibliography{local.bib}

% -----------------------------------------------------------------------------
% Fonts
% -----------------------------------------------------------------------------
\usepackage{../common/monster2e}

% -----------------------------------------------------------------------------
% Algorithms
% -----------------------------------------------------------------------------
\usepackage{../common/algorithm,../common/algorithmic,caption}

% -----------------------------------------------------------------------------
% Layout and positioning
% -----------------------------------------------------------------------------
\usepackage[absolute,overlay]{textpos}
\TPshowboxesfalse
\textblockorigin{0mm}{0mm}

% Graphics path
\graphicspath{{../img/}}

% -----------------------------------------------------------------------------
% Colors and symbols
% -----------------------------------------------------------------------------
% === MERGED misosymbols ===
% =============================================================================
% MVA Graphs in Machine Learning - Color Definitions
% =============================================================================

% -----------------------------------------------------------------------------
% 2025 Modern Color Palette
% Clean, cohesive palette for professional presentations
% -----------------------------------------------------------------------------

% Primary colors
\definecolor{mvaPrimary}{HTML}{1e3a5f}      % Deep navy - main brand color
\definecolor{mvaAccent}{HTML}{0891b2}       % Teal/cyan - accent and highlights
\definecolor{mvaHighlight}{HTML}{fbbf24}    % Amber - callouts and emphasis
\definecolor{mvaSuccess}{HTML}{059669}      % Emerald - positive/success states
\definecolor{mvaAlert}{HTML}{dc2626}        % Red - warnings and alerts

% Neutral colors
\definecolor{mvaBackground}{HTML}{f8fafc}   % Slate-50 - light backgrounds
\definecolor{mvaText}{HTML}{334155}         % Slate-700 - body text
\definecolor{mvaMuted}{HTML}{94a3b8}        % Slate-400 - muted/secondary text

% Additional modern colors
\definecolor{mvaLightAccent}{HTML}{67e8f9}  % Cyan-300 - light accent
\definecolor{mvaDarkAccent}{HTML}{0e7490}   % Cyan-700 - dark accent
\definecolor{mvaSubtle}{HTML}{e2e8f0}       % Slate-200 - subtle borders

% -----------------------------------------------------------------------------
% Legacy Color Definitions (kept for backward compatibility)
% -----------------------------------------------------------------------------
\definecolor{graphicbackground}{rgb}{0.96,0.96,0.8}
\definecolor{rouge1}{RGB}{226,0,38}         % red P
\definecolor{orange1}{RGB}{243,154,38}      % orange P
\definecolor{jaune}{RGB}{254,205,27}        % jaune P
\definecolor{blanc}{RGB}{255,255,255}       % blanc P
\definecolor{rouge2}{RGB}{230,68,57}        % red S
\definecolor{orange2}{RGB}{236,117,40}      % orange S
\definecolor{taupe}{RGB}{134,113,127}       % taupe S
\definecolor{gris}{RGB}{91,94,111}          % gris S
\definecolor{bleu1}{RGB}{38,109,131}        % bleu S
\definecolor{bleu2}{RGB}{28,50,114}         % bleu S
\definecolor{vert1}{RGB}{133,146,66}        % vert S
\definecolor{vert3}{RGB}{20,200,66}         % vert S
\definecolor{vert2}{RGB}{157,193,7}         % vert S
\definecolor{darkyellow}{RGB}{233,165,0}    % orange S
\definecolor{orangered}{RGB}{255,69,0}      % orange S
\definecolor{lightgray}{rgb}{0.9,0.9,0.9}
\definecolor{darkgray}{rgb}{0.666,0.666,0.666}
\definecolor{babyblue}{rgb}{0.54, 0.81, 0.94}
\definecolor{citrine}{rgb}{0.89, 0.82, 0.04}

% -----------------------------------------------------------------------------
% Highlight Color Variants (for equation/text highlighting)
% -----------------------------------------------------------------------------
\colorlet{redp}{mvaAlert!25}
\colorlet{greenp}{mvaSuccess!25}
\colorlet{bluep}{mvaAccent!25}
\colorlet{misogreen}{mvaSuccess!30}

% -----------------------------------------------------------------------------
% Text Color Commands
% Semantic coloring for emphasis and annotations
% -----------------------------------------------------------------------------

% Primary emphasis colors
\newcommand{\rcol}[1]{\textcolor{mvaAlert}{\textit{#1}}}
\newcommand{\defcol}[1]{\textcolor{mvaSuccess}{\textbf{#1}}}
\newcommand{\yellcol}[1]{\textcolor{mvaAccent}{\textbf{#1}}}
\newcommand{\gcol}[1]{\textcolor{mvaSuccess}{\textbf{#1}}}

\newcommand{\emphcol}[1]{\textcolor{mvaAccent!80!mvaText}{\textbf{#1}}}

% Muted/note colors
\newcommand{\grcol}[1]{\textcolor{mvaMuted}{#1}}
\newcommand{\notecol}[1]{\textcolor{mvaMuted}{#1}}

% Additional accent colors
\newcommand{\ycol}[1]{\textcolor{mvaHighlight!80!mvaText}{\textit{#1}}}

% Multi-color scheme for contrasting elements
\newcommand{\mcolA}[1]{{\color{mvaPrimary} #1}}
\newcommand{\mcolB}[1]{{\color{mvaHighlight!80!mvaText} #1}}
\newcommand{\mcolC}[1]{{\color{mvaSuccess} #1}}  


\DeclareMathOperator*{\argmax}{arg\,max}
\DeclareMathOperator*{\argmin}{arg\,min}
\DeclareMathOperator*{\arginf}{arg\,inf}
\DeclareMathOperator*{\sgn}{sgn}
%\DeclareMathOperator{\regret}{regret}
\DeclareMathOperator{\polylog}{polylog}
\DeclareMathOperator{\logloglog}{logloglog}
\DeclareMathOperator{\polyloglog}{polyloglog}


\newcommand*\diff{\mathop{}\!\mathrm{d}}
\newcommand*\Diff[1]{\mathop{}\!\mathrm{d^#1}}
\renewcommand{\d}[1]{\ensuremath{\operatorname{d}\!{#1}}}

\newcommand{\set}[1]{\left\{#1\right\}}
%\DeclarePairedDelimiter\ceil{\lceil}{\rceil}
%\DeclarePairedDelimiter\floor{\lfloor}{\rfloor}
\newcommand{\ceil}[1]{\left\lceil#1\right\rceil}
\newcommand{\floor}[1]{\left\lfloor#1\right\rfloor}

%\newcommand{\ceil}[1]{\lceil#1\rceil}

\newcommand{\II}[1]{\mathds{1}_{\left\{#1\right\}}}
\newcommand{\I}{{\mathds{1}}}


%arrows 
\newcommand{\ra}{\rightarrow}

%distributions 
\newcommand{\Bernoulli}{\mathrm{Bernoulli}}


\newcommand{\specialcell}[2][c]{%
 \begin{tabular}[#1]{@{}c@{}}#2\end{tabular}}

\newtheorem{assumption}{Assumption}
%\newtheorem{lemma}{Lemma}
%\newtheorem{theorem}{Theorem}
%\newtheorem{definition}{Definition}
%\newtheorem{corollary}{Corollary}
%\newtheorem{remark}{Remark}
\newtheorem{proposition}{Proposition}
\newtheorem{remark}{Remark}


% \newcommand{\R}{I\!\! R}
\newcommand{\R}{\mathbb{R}}
\newcommand{\realset}{\mathbb{R}}

\newcommand{\NN}{{\mathbb N}}
\newcommand{\1}{\mathds{1}}
\newcommand{\bOne}{{\bf 1}}
\newcommand{\bZero}{{\bf 0}}
\newcommand{\E}{\mathbb{E}}
\newcommand{\EE}[1]{\mathbb{E}\left[#1\right]}
\newcommand{\EEt}[1]{\mathbb{E}_t\left[#1\right]}
\newcommand{\EEs}[2]{\mathbb{E}_{#1}\left[#2\right]}
\newcommand{\EEc}[2]{\mathbb{E}\left[#1\left|#2\right.\right]}
\newcommand{\EEcc}[2]{\mathbb{E}\left[\left.#1\right|#2\right]}
\newcommand{\EEcct}[2]{\mathbb{E}_t\left[\left.#1\right|#2\right]}
\newcommand{\PP}[1]{\mathbb{P}\left[#1\right]}
\newcommand{\PPt}[1]{\mathbb{P}_t\left[#1\right]}
 \newcommand{\PPc}[2]{\mathbb{P}\left[#1\left|#2\right.\right]}
\newcommand{\PPcc}[2]{\mathbb{P}\left[\left.#1\right|#2\right]}
\newcommand{\PPct}[2]{\mathbb{P}_t\left[#1\left|#2\right.\right]}
\newcommand{\PPcct}[2]{\mathbb{P}_t\left[\left.#1\right|#2\right]}
%parens
\newcommand{\pa}[1]{\left(#1\right)}
\newcommand{\sqpa}[1]{\left[#1\right]}
\newcommand{\ac}[1]{\left\{#1\right\}}
\newcommand{\ev}[1]{\left\{#1\right\}}
\newcommand{\card}[1]{\left|#1\right|}


\newcommand{\normtwo}[1]{\|#1\|_2}
\newcommand{\norm}[1]{\left\|#1\right\|}
\newcommand{\onenorm}[1]{\norm{#1}_1}
\newcommand{\infnorm}[1]{\norm{#1}_\infty}
\newcommand{\norminf}[1]{\infnorm{#1}}

\newcommand{\abs}[1]{\left|#1\right|}

\newcommand{\CommaBin}{\mathbin{\raisebox{0.5ex}{,}}}
\newcommand*{\MyDef}{\mathrm{\tiny def}}
\newcommand*{\eqdefU}{\ensuremath{\mathop{\overset{\MyDef}{=}}}}% Unscaled version
\newcommand*{\eqdef}{\mathop{\overset{\MyDef}{\resizebox{\widthof{\eqdefU}}{\heightof{=}}{=}}}}
%\newcommand{\eqdef}{\stackrel{{\rm def}}{=}}
%\newcommand{\eqdef}{\stackrel{{\rm \tiny def}}{=}}
\newcommand{\transpose}{^\mathsf{\scriptscriptstyle T}}

%Calligraphic Shorthands
\newcommand{\cA}{\mathcal{A}}
\newcommand{\cB}{\mathcal{B}}
\newcommand{\cC}{\mathcal{C}}
\newcommand{\cD}{\mathcal{D}}
\newcommand{\cE}{\mathcal{E}}
\newcommand{\F}{\mathcal{F}}
\newcommand{\cF}{\mathcal{F}}
\newcommand{\cG}{\mathcal{G}}
\newcommand{\cH}{\mathcal{H}}
\newcommand{\cI}{\mathcal{I}}
\newcommand{\cJ}{\mathcal{J}}
\newcommand{\cK}{\mathcal{K}}
\newcommand{\cL}{\mathcal{L}}
\newcommand{\calL}{\cL}
\newcommand{\cM}{\mathcal{M}}
\newcommand{\cN}{\mathcal{N}}
\newcommand{\cO}{\mathcal{O}}
\newcommand{\tcO}{\widetilde{\cO}}
\newcommand{\OO}{\mathcal{O}}
\newcommand{\tOO}{\wt{\OO}}
\newcommand{\cP}{\mathcal{P}}
\newcommand{\cQ}{\mathcal{Q}}
\newcommand{\cR}{\mathcal{R}}
\newcommand{\Sw}{\mathcal{S}}
\newcommand{\cS}{\mathcal{S}}
\newcommand{\cT}{\mathcal{T}}
\newcommand{\T}{\cT}
\newcommand{\cU}{\mathcal{U}}
\newcommand{\cV}{\mathcal{V}}
\newcommand{\cW}{\mathcal{W}}
\newcommand{\cX}{\mathcal{X}}
\newcommand{\X}{\cX}
\newcommand{\cY}{\mathcal{Y}}
\newcommand{\cZ}{\mathcal{Z}}

%Bolds Shorthands
\newcommand{\ba}{{\bf a}}
\newcommand{\bA}{{\bf A}}
\newcommand{\bb}{{\bf b}}
\newcommand{\bB}{{\bf B}}
\newcommand{\bc}{{\bf c}}
\newcommand{\bC}{{\bf C}}
\newcommand{\bD}{{\bf D}}
\newcommand{\bg}{{\bf g}}
\newcommand{\bG}{{\bf G}}
\newcommand{\bI}{{\bf I}}
\newcommand{\bM}{{\bf M}}
\newcommand{\bO}{\boldsymbol{O}}
\newcommand{\bp}{\boldsymbol{p}}
\newcommand{\bP}{{\bf P}}
\newcommand{\br}{{\bf r}}
\newcommand{\bR}{{\bf R}}
\newcommand{\bQ}{{\bf Q}}
\newcommand{\be}{{\bf e}}
\newcommand{\bff}{{\bf f}}
\newcommand{\bi}{{\bf i}}
\newcommand{\bk}{{\bf k}}
\newcommand{\bK}{{\bf K}}
\newcommand{\bL}{{\bf L}}
\newcommand{\bs}{{\bf s}}
\newcommand{\bq}{{\bf q}}
\newcommand{\bu}{{\bf u}}
\newcommand{\bU}{{\bf U}}
\newcommand{\bv}{{\bf v}}
\newcommand{\bV}{{\bf V}}
\newcommand{\bw}{{\bf w}}
\newcommand{\bW}{{\bf W}}
\newcommand{\by}{{\bf y}}
\newcommand{\bx}{{\bf x}}
\newcommand{\bX}{{\bf X}}
\newcommand{\bZ}{{\bf Z}}

\newcommand{\eps}{\varepsilon}
\renewcommand{\epsilon}{\varepsilon}
\renewcommand{\hat}{\widehat}
\renewcommand{\tilde}{\widetilde}
\renewcommand{\bar}{\overline}

\newcommand{\balpha}{{\boldsymbol \alpha}}
\newcommand{\talpha}{\widetilde{\alpha}}
\newcommand{\btheta}{{\boldsymbol \theta}}
\newcommand{\tTheta}{{\widetilde\Theta}}
\newcommand{\bdelta}{{\boldsymbol \delta}}
\newcommand{\bDelta}{{\boldsymbol \Delta}}
\newcommand{\bLambda}{{\boldsymbol \Lambda}}
\newcommand{\bSigma}{{\boldsymbol \Sigma}}
\newcommand{\bmu}{{\boldsymbol \mu}}
\newcommand{\bxi}{{\boldsymbol \xi}}
\newcommand{\bell}{\boldsymbol \ell}

\newcommand{\nothere}[1]{}
\newcommand{\moveb}{\\ \bigskip}
\newcommand{\movebb}{\\\[-0.25em]}
\newcommand{\moves}{\\ \smallskip}
\newcommand{\movess}{\\\[-0.5em]}
\newcommand{\movesss}{\smallskip}
\newcommand{\bskip}{\vspace{\baselineskip}}
\newcommand{\hbskip}{\vspace{.5\baselineskip}}

% bandits
\newcommand{\hloss}{\hat\ell}
\newcommand{\bloss}{\boldsymbol  \ell}
\newcommand{\hbl}{\hat{\bloss}}
\newcommand{\hbL}{\wh{\bL}}
\newcommand{\wh}{\widehat}
\newcommand{\ti}{_{t,i}}
\newcommand{\wt}{\widetilde}


%% from single papers, but merged here
% algo names
\usepackage{xspace}
\renewcommand{\ttdefault}{lmtt}
%other algos
\newcommand{\LP}{\texttt{LP}\xspace}
\newcommand{\CMG}{\texttt{CMG}\xspace}
%bandits
\newcommand{\FPL}{\texttt{FPL}\xspace}
\newcommand{\TS}{\texttt{TS}\xspace}
\newcommand{\UCB}{\texttt{UCB}\xspace}
\newcommand{\MOSS}{\texttt{MOSS}\xspace}
\newcommand{\UCBE}{\texttt{UCB-E}\xspace}
\newcommand{\ImprovedUCB}{\texttt{ImprovedUCB}\xspace}
\newcommand{\klucb}{\texttt{KL-UCB}\xspace}
\newcommand{\CUCB}{\texttt{CUCB}\xspace}
\newcommand{\EXP}{\texttt{Exp3}\xspace}
\newcommand{\exph}{\EXP}
%linear bandits
\newcommand{\LinearTS}{\texttt{LinearTS}\xspace}
\newcommand{\ThompsonSampling}{\texttt{ThompsonSampling}\xspace}
\newcommand{\SpectralEliminator}{\texttt{\textcolor[rgb]{0.5,0.2,0}{SpectralEliminator}}\xspace}
\newcommand{\LinearEliminator}{\texttt{\textcolor[rgb]{0.5,0.2,0}{LinearEliminator}}\xspace}
\newcommand{\LinUCB}{\texttt{LinUCB}\xspace}
\newcommand{\LinRel}{\texttt{LinRel}\xspace}
\newcommand{\KernelUCB}{\texttt{\textcolor[rgb]{0.5,0.2,0}{KernelUCB}}\xspace}
\newcommand{\SupKernelUCB}{\texttt{\textcolor[rgb]{0.5,0.2,0}{SupKernelUCB}}\xspace}
\newcommand{\GPUCB}{\texttt{GP-UCB}\xspace}
\newcommand{\OFUL}{\texttt{OFUL}\xspace}
\newcommand{\OPM}{\texttt{\textcolor[rgb]{0.5,0.2,0}{OPM}}\xspace}
%graph bandits
\newcommand{\CLUB}{\texttt{CLUB}\xspace}
\newcommand{\GOBLin}{\texttt{GOB.Lin}\xspace}
\newcommand{\UCBN}{\texttt{UCB-N}\xspace}
\newcommand{\UCBmaxN}{\texttt{UCB-MaxN}\xspace}
\newcommand{\GraphMOSS}{\texttt{\textcolor[rgb]{0.5,0.2,0}{GraphMOSS}}\xspace}
\newcommand{\SpectralUCB}{\texttt{\textcolor[rgb]{0.5,0.2,0}{SpectralUCB}}\xspace}
\newcommand{\CheapUCB}{\texttt{\textcolor[rgb]{0.5,0.2,0}{CheapUCB}}\xspace}
\newcommand{\SpectralTS}{\texttt{\textcolor[rgb]{0.5,0.2,0}{SpectralTS}}\xspace}
\newcommand{\SupLinRel}{\texttt{SupLinRel}\xspace}
\newcommand{\SupLinUCB}{\texttt{SupLinUCB}\xspace}
\newcommand{\imb}{\texttt{\textcolor[rgb]{0.5,0.2,0}{IMLinUCB}}\xspace}
\newcommand{\NetBandits}{\texttt{NetBandits}\xspace}
\newcommand{\BARE}{\texttt{\textcolor[rgb]{0.5,0.2,0}{BARE}}\xspace}
\newcommand{\ELP}{\texttt{ELP}\xspace}
\newcommand{\ELPP}{\texttt{ELP.P}\xspace}
\newcommand{\expix}{\texttt{\textcolor[rgb]{0.5,0.2,0}{Exp3-IX}}\xspace}
\newcommand{\expset}{\texttt{Exp3-SET}\xspace}
\newcommand{\expdom}{\texttt{Exp3-DOM}\xspace}
\newcommand{\expg}{\texttt{Exp3.G}\xspace}
\newcommand{\fplbgr}{\texttt{FPL-BGR}\xspace}
\newcommand{\fplix}{\texttt{\textcolor[rgb]{0.5,0.2,0}{FPL-IX}}\xspace}
\newcommand{\comphedge}{\texttt{Component\-Hedge}\xspace}
\newcommand{\hedge}{\texttt{Hedge}\xspace}
\newcommand{\expxxx}{\texttt{\textcolor[rgb]{0.5,0.2,0}{Exp3-WIX}}\xspace}
\newcommand{\expwix}{\texttt{\textcolor[rgb]{0.5,0.2,0}{Exp3-WIX}}\xspace}
\newcommand{\expixa}{\texttt{Exp3-IXa}\xspace}
\newcommand{\expixb}{\texttt{Exp3-IXb}\xspace}
\newcommand{\expixt}{\texttt{Exp3-IXt}\xspace}
\newcommand{\expcoop}{\texttt{Exp3-Coop}\xspace}
\newcommand{\expres}{\texttt{\textcolor[rgb]{0.5,0.2,0}{Exp3-Res}}\xspace}

%continuous bandits
\newcommand{\StoSOO}{\texttt{\textcolor[rgb]{0.5,0.2,0}{StoSOO}}\xspace}
\newcommand{\POO}{\texttt{\textcolor[rgb]{0.5,0.2,0}{POO}}\xspace}
\newcommand{\DOO}{\texttt{DOO}\xspace}
\newcommand{\SOO}{\texttt{SOO}\xspace}
\newcommand{\Zooming}{\texttt{Zooming}\xspace}
\newcommand{\UCT}{\texttt{UCT}\xspace}
\newcommand{\HCT}{\texttt{HCT}\xspace}
\newcommand{\SHOO}{\POO}
\newcommand{\HOO}{\texttt{HOO}\xspace}
\newcommand{\ATB}{\texttt{ATB}\xspace}
\newcommand{\TZ}{\texttt{TaxonomyZoom}\xspace}
\newcommand{\Direct}{\texttt{DiRect}\xspace}
%extreme bandits
\newcommand{\SiRI}{\texttt{\textcolor[rgb]{0.5,0.2,0}{SiRI}}\xspace}
%planning "bandits"
\newcommand{\olop}{\texttt{OLOP}\xspace}
\newcommand{\stopalgo}{\texttt{StOP}\xspace}
\newcommand{\metagrill}{\texttt{\textcolor[rgb]{0.5,0.2,0}{\textup{TrailBlazer}}}\xspace}
%polymatroid bandits
\newcommand{\greedy}{\texttt{Greedy}\xspace}
\newcommand{\opm}{\texttt{\textcolor[rgb]{0.5,0.2,0}{OPM}}\xspace}

%other algorithms
\newcommand{\pagerank}{\texttt{PageRank}\xspace}


\newcommand{\maxn}{\texttt{max}\xspace}
\newcommand{\avgn}{\texttt{avg}\xspace}


%kernelUCB
\newcommand{\reg}{\gamma}
\newcommand{\hmu}{\hat{\mu}}
%\newcommand{\hm}{\hat{\mu}}
\newcommand{\hw}{\hat{w}}
\newcommand{\hth}{\hat{\theta}}
\newcommand{\hs}{\hat{\sigma}}
\newcommand{\td}{\tilde{d}}

%graph bandit notation
\newcommand{\etat}{\eta_t}
\newcommand{\gammat}{\gamma_t}
\newcommand{\nodes}{{\textcolor{green!80!black}{N}}}
\newcommand{\rounds}{{\textcolor[rgb]{0.25,0.0,0.6}{T}}}
\newcommand{\nodeset}{\cV}
\newcommand{\edgeset}{\cE}
\newcommand{\regret}{R_\rounds}
\newcommand{\cgamma}{c_\gamma}
\newcommand{\cgammat}{c_{\gamma_t}}
\newcommand{\sumT}{\sum_{t = 1}^\rounds}
\newcommand{\sumt}{\sum_{t=1}^\rounds}
\newcommand{\sumtl}{\sum\limits_{t=1}^\rounds}
\newcommand{\sumj}{\sum_{j\in \nodes_i^-}}
\newcommand{\sumtj}{\sum_{j\in \nodes_{t,i}^-}}
\newcommand{\sumi}{\sum_{i=1}^{\nodes}}
\newcommand{\sumji}{\sum_{j\in \{\nodes_i^-\cup\{i\}\}}}
\newcommand{\sumtji}{\sum_{j\in \{\nodes_{t,i}^-\cup\{i\}\}}}
\newcommand{\dti}{d_{t,i}^-}
\newcommand{\hdi}{\hat{d}_i^-}
\newcommand{\hdk}{\hat{d}_k^-}
\newcommand{\hd}{\hat{d}^-}
\newcommand{\tti}{_{t+1,i}}
\newcommand{\tj}{_{t,j}}
\newcommand{\ji}{_{j,i}}
\newcommand{\Ii}{_{I_t,i}}
\newcommand{\pti}{p\ti}
\newcommand{\pta}{p_{t,a}}
\newcommand{\qti}{q\ti}
\newcommand{\hpti}{\hat{p}\ti}
\newcommand{\hpi}{\hat{p}_i}
\newcommand{\hp}{\hat{p}}
\newcommand{\hqti}{\hat{q}\ti}
\newcommand{\ptj}{p_{t,j}}
\newcommand{\qtj}{q_{t,j}}
\newcommand{\hptj}{\hat{p}_{t,j}}
\newcommand{\hqtj}{\hat{q}_{t,j}}
\newcommand{\oti}{o\ti}
\newcommand{\Oti}{O\ti}
\newcommand{\loss}{\ell}
\newcommand{\hLoss}{\hat{L}}
\newcommand{\hL}{\wh{L}}
\newcommand{\noise}{\xi}
%spectral
\newcommand{\dold}{d_{\scriptsize\mbox{old}}}
\newcommand{\dnew}{d_{\scriptsize\mbox{new}}}
%wix
\newcommand{\gweight}{s}
\newcommand{\avgalpha}{\alpha^*_{\text{avg}}}
%bare
\newcommand{\rkdual}{r_k^{\circ}}
\newcommand{\rktdual}{r_{k,t}^{\circ}}
\newcommand{\rkprimetdual}{r_{k',t}^{\circ}}
\newcommand{\rkttdual}{r_{k,t+1}^{\circ}}
\newcommand{\rkprimettdual}{r_{k',t+1}^{\circ}}
\newcommand{\ridual}{r_i^{\circ}}
\newcommand{\rstardual}{r_\star^{\circ}}
\newcommand{\Ddual}{D^{\circ}}
\newcommand{\Ddualset}{\mathcal D^{\circ}}
%continuous bandit notation
%\newcommand{\node}[2]{\circ[#1,#2]}
\newcommand{\node}[2]{(#1,#2)}


% =============================================================================
% From mathdef.tex (Paper Dependent Stuff)
% =============================================================================

\newcommand{\disre}{\textsc{DisRe}\xspace}
\newcommand{\hfssparse}{\textsc{Sparse-HFS}\xspace}
\newcommand{\hfsstable}{\textsc{Stable-HFS}\xspace}
\newcommand{\hfs}{\texttt{HFS}\xspace}
\newcommand{\rls}{\texttt{LapSmo}\xspace}
\newcommand{\stahfs}{\texttt{STA}\xspace}
\newcommand{\simplehfs}{\textsc{Simple-HFS}\xspace}
\newcommand{\subs}{\textsc{SubSampling}\xspace}
\newcommand{\fergus}{\textsc{EigFun}\xspace}
\newcommand{\onenn}{\textsc{1-NN}\xspace}

\newcommand{\qbar}{\wb{q}}
\newcommand{\ssl}{\textsc{SSL}\xspace}
\newcommand{\scluster}{\textsc{SC}\xspace}
\newcommand{\squeak}{\texttt{SQUEAK}\xspace}
\newcommand{\gd}{\texttt{GD}\xspace}
\newcommand{\lembed}{\textsc{LE}\xspace}
\newcommand{\kn}{\texttt{kN}\xspace}
\newcommand{\spect}{\textsc{spect}\xspace}
\newcommand{\gspect}{$\gamma569Xl\textsc{spect}\xspace}
\newcommand{\exact}{\texttt{EXACT}\xspace}
\newcommand{\ltr}{\texttt{LTR}\xspace}

\newcommand{\sampdist}{\rho}
\newcommand{\bphi}{\boldsymbol{\phi}}

\newcommand{\bomega}{\boldsymbol{\omega}}
\newcommand{\bupsilon}{\boldsymbol{\upsilon}}
%\newcommand{\bSigma}{\boldsymbol{\Sigma}}

\newcommand{\bwavec}{\bomega}
\newcommand{\buavec}{\boldsymbol{\upsilon}}
\newcommand{\bvavec}{\boldsymbol{\theta}}
\newcommand{\nystrom}{Nystr\"{o}m\xspace}
\newcommand{\feasibleset}{\mathcal{S}}
\newcommand{\coldict}{\mathcal{I}}
\newcommand{\dictentry}{\wb{\bx}}
\newcommand{\pvec}{\mathbf{v}}
\newcommand{\afeatvec}{\widetilde{\boldsymbol{\phi}}}
\newcommand{\pmat}{\mathbf{V}}

\newcommand{\onlinesqueak}{\textsc{KORS}\xspace}
\newcommand{\kons}{\textsc{KONS}\xspace}
\newcommand{\ons}{\textsc{ONS}\xspace}
\newcommand{\conskons}{\textsc{Con-KONS}\xspace}
\newcommand{\bskons}{\textsc{B-KONS}\xspace}
\newcommand{\batch}{\textsc{BATCH}\xspace}
\newcommand{\skons}{\textsc{PROS-N-KONS}\xspace}
\newcommand{\icmlskons}{\textsc{Sketched-KONS}\xspace}
\newcommand{\fixskons}{\textsc{Fixed-SKONS}\xspace}
\newcommand{\nysgd}{\textsc{NOGD}\xspace}
\newcommand{\fogd}{\textsc{FOGD}\xspace}
\newcommand{\dualgd}{\textsc{Dual-SGD}\xspace}

% Aesthetics
\def\:#1{\protect \ifmmode {\mathbf{#1}} \else {\textbf{#1}} \fi}
\newcommand{\wrt}{w.r.t.\xspace}
\newcommand{\defeq}{\stackrel{\mathclap{\normalfont\mbox{\tiny def}}}{=}}
\newcommand{\maxund}[1]{\max\limits_{#1}}
\newcommand{\minund}[1]{\min\limits_{#1}}
\newcommand{\bigotime}{\mathcal{O}}
\newcommand{\bigomegatime}{{\Omega}}
\DeclareMathOperator*{\poly}{poly}

\newcommand{\normempty}[1]{\left\Vert #1 \right\Vert}
\DeclareMathOperator*{\Op}{\textbf{op}}
\newcommand{\transp}{\mathsf{\scriptscriptstyle T}}
\DeclareMathOperator*{\Tr}{Tr}
\DeclareMathOperator*{\Det}{Det}
\DeclareMathOperator*{\Ran}{Im}
\DeclareMathOperator*{\Ker}{Ker}
\DeclareMathOperator*{\Diag}{Diag}
\DeclareMathOperator*{\Vecm}{Vec}
\DeclareMathOperator*{\Rank}{Rank}
\DeclareMathOperator*{\Dim}{Dim}
\DeclareMathOperator*{\Row}{Row}
\DeclareMathOperator*{\Col}{Col}
\DeclareMathOperator*{\Nnz}{nnz}

\newcommand{\probability}{\mathbb{P}}
\newcommand{\probdist}{Pr}
\DeclareMathOperator*{\expectedvalue}{\mathbb{E}}
\DeclareMathOperator*{\variance}{\text{Var}}
\newcommand{\expectedvalueover}[1]{\expectedvalue\limits_{#1}}
\newcommand{\condbar}{\;\middle|\;}
\newcommand{\gaussdistr}{\mathcal{N}}
\newcommand{\uniformdistr}{\mathcal{U}}
\newcommand{\bernoullidist}{\mathcal{B}}
\newcommand{\filtration}{\mathcal{F}}
\newcommand{\learningalg}{\mathcal{L}}
\newcommand{\indvec}{J}
\newcommand{\indfunc}{\mathbb{I}}

\newcommand{\Real}{\mathbb{R}}
\newcommand{\Natural}{\mathbb{N}}
\newcommand{\statespace}{\mathcal{X}}
\newcommand{\funcspace}{\mathcal{F}}
\newcommand{\dataset}{\mathcal{D}}
\newcommand{\trainingset}{\mathcal{S}}
\newcommand{\testset}{\mathcal{T}}

\newcommand{\laplacian}{\bL}
\newcommand{\adjacency}{A}
\newcommand{\Gg}{\mathcal{G}}
\newcommand{\Hg}{\mathcal{H}}
\newcommand{\Ag}{\mathcal{A}}
\newcommand{\vertexset}{\mathcal{V}}

\newcommand{\rkhs}{\mathcal{H}}
\newcommand{\func}{f}
\newcommand{\whfunc}{\wh{\mathsf{f}}}
\newcommand{\wtfunc}{\wt{\mathsf{f}}}
\newcommand{\hilbprod}[2]{\left\langle{#1},{#2}\right\rangle_\rkhs}

\newcommand{\kerfunc}{\mathcal{K}}
\newcommand{\kermatrix}{{\bK}}
\newcommand{\featkermatrix}{\boldsymbol{\Phi}}
\newcommand{\featvec}{\boldsymbol{\phi}}
\newcommand{\deff}{d_{\text{eff}}}
\newcommand{\donl}{d_{\text{onl}}}
\newcommand{\bdeff}[1]{\wb{d}_{\text{eff}}(#1)}
\newcommand{\adeff}[1]{\wt{d}_{\text{eff}}(#1)}
\newcommand{\atau}{\wt{\tau}}
\newcommand{\axi}{\wt{\xi}}
\newcommand{\akermatrix}{\wt{\mathbf{K}}}
\newcommand{\bkermatrix}{\mathbf{\wb{K}}}
\newcommand{\selmatrix}{{\bS}}
\newcommand{\selrmatrix}{{\bR}}

% End of From mathdef.tex

% === MERGED highlights ===
% =============================================================================
% MVA Graphs in Machine Learning - Shared Highlight Commands
% Unified text and equation highlighting for all subdecks
% =============================================================================

% -----------------------------------------------------------------------------
% Required packages (should be loaded in misomva.tex)
% -----------------------------------------------------------------------------
\usepackage[most]{tcolorbox}

% -----------------------------------------------------------------------------
% Equation Highlight Commands
% Colored background for inline math
% -----------------------------------------------------------------------------
\newcommand{\highlight}[2]{%
  \colorbox{#1}{$\displaystyle#2$}}

% Equation highlights with semantic colors
\newcommand{\hlrEq}[1]{\highlight{redp}{#1}}      % Red highlight for equations
\newcommand{\hlbEq}[1]{\highlight{bluep}{#1}}     % Blue highlight for equations
\newcommand{\hlgEq}[1]{\highlight{greenp}{#1}}    % Green highlight for equations

% -----------------------------------------------------------------------------
% Text Highlight Commands
% Colored background for text
% -----------------------------------------------------------------------------
\newcommand{\hl}[3][\fboxsep1pt]{{#1\colorbox{#2}{#3}}}

% Text highlights with semantic colors
\newcommand{\hlr}[1]{\hl{redp}{#1}}               % Red highlight
\newcommand{\hlg}[1]{\hl{greenp}{#1}}             % Green highlight
\newcommand{\hlb}[1]{\hl{bluep}{#1}}              % Blue highlight

% Bold text highlights
\newcommand{\hlbr}[1]{\hl[\fboxsep0pt]{redp}{$\mathbf{#1}$}}
\newcommand{\hlbg}[1]{\hl[\fboxsep0pt]{greenp}{$\mathbf{#1}$}}
\newcommand{\hlbb}[1]{\hl[\fboxsep0pt]{bluep}{$\mathbf{#1}$}}

% -----------------------------------------------------------------------------
% Block Quote Environment
% Modern styled callout boxes for important content
% -----------------------------------------------------------------------------
\newtcolorbox{blockquote}{
  colback=mvaAccent!10,
  colframe=mvaAccent!30,
  grow to right by=-1mm,
  grow to left by=-1mm,
  boxrule=1pt,
  boxsep=0pt,
  left=8pt,
  right=8pt,
  top=6pt,
  bottom=6pt,
  breakable,
  arc=4pt
}

% Alternative callout styles
\newtcolorbox{infobox}{
  colback=mvaPrimary!8,
  colframe=mvaPrimary!40,
  boxrule=1pt,
  boxsep=0pt,
  left=8pt,
  right=8pt,
  top=6pt,
  bottom=6pt,
  arc=4pt
}

\newtcolorbox{alertbox}{
  colback=mvaAlert!8,
  colframe=mvaAlert!40,
  boxrule=1pt,
  boxsep=0pt,
  left=8pt,
  right=8pt,
  top=6pt,
  bottom=6pt,
  arc=4pt
}

\newtcolorbox{successbox}{
  colback=mvaSuccess!8,
  colframe=mvaSuccess!40,
  boxrule=1pt,
  boxsep=0pt,
  left=8pt,
  right=8pt,
  top=6pt,
  bottom=6pt,
  arc=4pt
}

% =============================================================================




% -----------------------------------------------------------------------------
% Shared TikZ styles for graph visualization
% -----------------------------------------------------------------------------
% === MERGED tikzstyles ===
% =============================================================================
% MVA Graphs in Machine Learning - Shared TikZ Styles
% Unified graph visualization styles for all subdecks
% =============================================================================

% -----------------------------------------------------------------------------
% Graph Vertex Styles
% -----------------------------------------------------------------------------
\tikzset{
  % Standard vertex style
  vertex/.style={
    circle,
    fill=mvaPrimary!20,
    minimum size=20pt,
    inner sep=0pt,
    font=\small
  },
  % Alternative name for compatibility
  mvaVertex/.style={vertex},
  % Selected/highlighted vertex
  selected vertex/.style={
    vertex,
    fill=mvaAccent!40
  },
  mvaHighlightVertex/.style={selected vertex},
}

% -----------------------------------------------------------------------------
% Graph Edge Styles
% -----------------------------------------------------------------------------
\tikzset{
  % Standard edge style
  edge/.style={
    draw,
    thick,
    ->,
    mvaPrimary!80
  },
  mvaEdge/.style={edge},
  % Edge weight label
  weight/.style={
    font=\small
  },
  % Selected/highlighted edge
  selected edge/.style={
    draw,
    line width=5pt,
    -,
    mvaAccent!50
  },
  mvaHighlightEdge/.style={selected edge},
  % Ignored/dimmed edge
  ignored edge/.style={
    draw,
    line width=5pt,
    -,
    mvaMuted!30
  }
}

% -----------------------------------------------------------------------------
% Special Node Styles
% -----------------------------------------------------------------------------
\tikzset{
  % Labeled node (positive class)
  positive node/.style={
    vertex,
    fill=mvaSuccess!40
  },
  % Labeled node (negative class)
  negative node/.style={
    vertex,
    fill=mvaAlert!40
  },
  % Unlabeled node
  unlabeled node/.style={
    vertex,
    fill=mvaMuted!20
  }
}

% -----------------------------------------------------------------------------
% Isara Logo - TikZ Vector Graphic
% Beautiful vector representation of the Isara logo as a piece of art
% -----------------------------------------------------------------------------

% Helper commands to draw individual logo elements (for artistic use)
% Accurate TikZ recreation of the Isara logo based on actual logo design
\newcommand{\drawIsaraCircle}[1]{% #1 = color - double concentric circles
  \draw[#1, line width=0.08, line cap=round] (0,0) circle (1.0);
  \draw[#1, line width=0.08, line cap=round] (0,0) circle (0.88);
}
\newcommand{\drawIsaraMainCurve}[1]{% #1 = color - central thick calligraphic flowing B-like shape
  % This is the dominant feature - a thick, smooth calligraphic stroke
  % Upper left loop
  \fill[#1]
    (-0.48, 0.28)
    .. controls (-0.62, 0.38) and (-0.68, 0.58) .. (-0.60, 0.72)
    .. controls (-0.52, 0.82) and (-0.35, 0.85) .. (-0.18, 0.82)
    .. controls (0.0, 0.78) and (0.15, 0.68) .. (0.25, 0.52)
    .. controls (0.32, 0.40) and (0.35, 0.25) .. (0.32, 0.10)
    .. controls (0.30, 0.0) and (0.25, -0.08) .. (0.18, -0.08)
    .. controls (0.12, -0.08) and (0.08, -0.02) .. (0.05, 0.08)
    % Central vertical segment
    .. controls (0.02, 0.18) and (-0.02, 0.28) .. (-0.10, 0.35)
    .. controls (-0.18, 0.42) and (-0.28, 0.45) .. (-0.38, 0.40)
    .. controls (-0.42, 0.38) and (-0.45, 0.34) .. (-0.48, 0.28)
    -- cycle;
  % Lower right loop
  \fill[#1]
    (0.48, -0.28)
    .. controls (0.62, -0.38) and (0.68, -0.58) .. (0.60, -0.72)
    .. controls (0.52, -0.82) and (0.35, -0.85) .. (0.18, -0.82)
    .. controls (0.0, -0.78) and (-0.15, -0.68) .. (-0.25, -0.52)
    .. controls (-0.32, -0.40) and (-0.35, -0.25) .. (-0.32, -0.10)
    .. controls (-0.30, 0.0) and (-0.25, 0.08) .. (-0.18, 0.08)
    .. controls (-0.12, 0.08) and (-0.08, 0.02) .. (-0.05, -0.08)
    % Central vertical segment
    .. controls (-0.02, -0.18) and (0.02, -0.28) .. (0.10, -0.35)
    .. controls (0.18, -0.42) and (0.28, -0.45) .. (0.38, -0.40)
    .. controls (0.42, -0.38) and (0.45, -0.34) .. (0.48, -0.28)
    -- cycle;
  % Central connecting stroke
  \fill[#1]
    (-0.08, 0.35)
    .. controls (-0.05, 0.30) and (0.0, 0.20) .. (0.0, 0.0)
    .. controls (0.0, -0.20) and (-0.05, -0.30) .. (-0.08, -0.35)
    .. controls (-0.12, -0.35) and (-0.18, -0.32) .. (-0.22, -0.25)
    .. controls (-0.18, -0.18) and (-0.15, 0.0) .. (-0.15, 0.0)
    .. controls (-0.15, 0.0) and (-0.18, 0.18) .. (-0.22, 0.25)
    .. controls (-0.18, 0.32) and (-0.12, 0.35) .. (-0.08, 0.35)
    -- cycle;
}
\newcommand{\drawIsaraTopSwirl}[1]{% #1 = color - thin outline curve upper left
  \draw[#1, line width=0.06, line cap=round]
    (-0.70, 0.35)
    .. controls (-0.72, 0.50) and (-0.68, 0.65) .. (-0.55, 0.75)
    .. controls (-0.40, 0.82) and (-0.22, 0.80) .. (-0.08, 0.72);
}
\newcommand{\drawIsaraBottomSwirl}[1]{% #1 = color - thin outline curve lower right
  \draw[#1, line width=0.06, line cap=round]
    (0.70, -0.35)
    .. controls (0.72, -0.50) and (0.68, -0.65) .. (0.55, -0.75)
    .. controls (0.40, -0.82) and (0.22, -0.80) .. (0.08, -0.72);
}
\newcommand{\drawIsaraLeftCurve}[1]{% #1 = color - thin outline curve right side of upper loop
  \draw[#1, line width=0.06, line cap=round]
    (0.38, 0.45)
    .. controls (0.48, 0.58) and (0.52, 0.72) .. (0.48, 0.82)
    .. controls (0.42, 0.88) and (0.32, 0.90) .. (0.20, 0.88);
}
\newcommand{\drawIsaraRightCurve}[1]{% #1 = color - thin outline curve left side of lower loop
  \draw[#1, line width=0.06, line cap=round]
    (-0.38, -0.45)
    .. controls (-0.48, -0.58) and (-0.52, -0.72) .. (-0.48, -0.82)
    .. controls (-0.42, -0.88) and (-0.32, -0.90) .. (-0.20, -0.88);
}
\newcommand{\drawIsaraInnerTopSwirl}[1]{% #1 = color - thin outline inner upper loop
  \draw[#1, line width=0.06, line cap=round]
    (-0.35, 0.52)
    .. controls (-0.25, 0.62) and (-0.10, 0.65) .. (0.05, 0.58)
    .. controls (0.15, 0.52) and (0.22, 0.42) .. (0.25, 0.28);
}
\newcommand{\drawIsaraInnerBottomSwirl}[1]{% #1 = color - thin outline inner lower loop
  \draw[#1, line width=0.06, line cap=round]
    (0.35, -0.52)
    .. controls (0.25, -0.62) and (0.10, -0.65) .. (-0.05, -0.58)
    .. controls (-0.15, -0.52) and (-0.22, -0.42) .. (-0.25, -0.28);
}
\newcommand{\drawIsaraInnerSwirl}[1]{% #1 = color, draws both inner swirls
  \drawIsaraInnerTopSwirl{#1}
  \drawIsaraInnerBottomSwirl{#1}
}

% Helper command to draw the complete Isara logo shape with gradient (used internally)
% This draws all elements with white-to-blue gradient matching the actual logo
\newcommand{\drawIsaraLogoShape}[1]{% #1 = base color (used for gradient creation)
  \def\logoRadius{1.0}
  % Define gradient shading from white to blue
  \shadedraw[top color=white, bottom color=mvaPrimary!60, line width=0.08, line cap=round]
    (0,0) circle (1.0);
  \shadedraw[top color=white, bottom color=mvaPrimary!60, line width=0.08, line cap=round]
    (0,0) circle (0.88);

  % Central thick calligraphic stroke with gradient
  % Upper left loop
  \shade[top color=white, bottom color=mvaPrimary!60]
    (-0.48, 0.28)
    .. controls (-0.62, 0.38) and (-0.68, 0.58) .. (-0.60, 0.72)
    .. controls (-0.52, 0.82) and (-0.35, 0.85) .. (-0.18, 0.82)
    .. controls (0.0, 0.78) and (0.15, 0.68) .. (0.25, 0.52)
    .. controls (0.32, 0.40) and (0.35, 0.25) .. (0.32, 0.10)
    .. controls (0.30, 0.0) and (0.25, -0.08) .. (0.18, -0.08)
    .. controls (0.12, -0.08) and (0.08, -0.02) .. (0.05, 0.08)
    .. controls (0.02, 0.18) and (-0.02, 0.28) .. (-0.10, 0.35)
    .. controls (-0.18, 0.42) and (-0.28, 0.45) .. (-0.38, 0.40)
    .. controls (-0.42, 0.38) and (-0.45, 0.34) .. (-0.48, 0.28)
    -- cycle;
  % Lower right loop
  \shade[top color=mvaPrimary!60, bottom color=white]
    (0.48, -0.28)
    .. controls (0.62, -0.38) and (0.68, -0.58) .. (0.60, -0.72)
    .. controls (0.52, -0.82) and (0.35, -0.85) .. (0.18, -0.82)
    .. controls (0.0, -0.78) and (-0.15, -0.68) .. (-0.25, -0.52)
    .. controls (-0.32, -0.40) and (-0.35, -0.25) .. (-0.32, -0.10)
    .. controls (-0.30, 0.0) and (-0.25, 0.08) .. (-0.18, 0.08)
    .. controls (-0.12, 0.08) and (-0.08, 0.02) .. (-0.05, -0.08)
    .. controls (-0.02, -0.18) and (0.02, -0.28) .. (0.10, -0.35)
    .. controls (0.18, -0.42) and (0.28, -0.45) .. (0.38, -0.40)
    .. controls (0.42, -0.38) and (0.45, -0.34) .. (0.48, -0.28)
    -- cycle;
  % Central connecting stroke
  \shade[left color=white, right color=mvaPrimary!50]
    (-0.08, 0.35)
    .. controls (-0.05, 0.30) and (0.0, 0.20) .. (0.0, 0.0)
    .. controls (0.0, -0.20) and (-0.05, -0.30) .. (-0.08, -0.35)
    .. controls (-0.12, -0.35) and (-0.18, -0.32) .. (-0.22, -0.25)
    .. controls (-0.18, -0.18) and (-0.15, 0.0) .. (-0.15, 0.0)
    .. controls (-0.15, 0.0) and (-0.18, 0.18) .. (-0.22, 0.25)
    .. controls (-0.18, 0.32) and (-0.12, 0.35) .. (-0.08, 0.35)
    -- cycle;

  % Outline curves with gradient
  \draw[white!90!mvaPrimary, line width=0.06, line cap=round]
    (-0.70, 0.35)
    .. controls (-0.72, 0.50) and (-0.68, 0.65) .. (-0.55, 0.75)
    .. controls (-0.40, 0.82) and (-0.22, 0.80) .. (-0.08, 0.72);
  \draw[white!90!mvaPrimary, line width=0.06, line cap=round]
    (0.70, -0.35)
    .. controls (0.72, -0.50) and (0.68, -0.65) .. (0.55, -0.75)
    .. controls (0.40, -0.82) and (0.22, -0.80) .. (0.08, -0.72);
  \draw[white!90!mvaPrimary, line width=0.06, line cap=round]
    (0.38, 0.45)
    .. controls (0.48, 0.58) and (0.52, 0.72) .. (0.48, 0.82)
    .. controls (0.42, 0.88) and (0.32, 0.90) .. (0.20, 0.88);
  \draw[white!90!mvaPrimary, line width=0.06, line cap=round]
    (-0.38, -0.45)
    .. controls (-0.48, -0.58) and (-0.52, -0.72) .. (-0.48, -0.82)
    .. controls (-0.42, -0.88) and (-0.32, -0.90) .. (-0.20, -0.88);
  \draw[white!90!mvaPrimary, line width=0.06, line cap=round]
    (-0.35, 0.52)
    .. controls (-0.25, 0.62) and (-0.10, 0.65) .. (0.05, 0.58)
    .. controls (0.15, 0.52) and (0.22, 0.42) .. (0.25, 0.28);
  \draw[white!90!mvaPrimary, line width=0.06, line cap=round]
    (0.35, -0.52)
    .. controls (0.25, -0.62) and (0.10, -0.65) .. (-0.05, -0.58)
    .. controls (-0.15, -0.52) and (-0.22, -0.42) .. (-0.25, -0.28);
}

% Main Isara logo command - draws the complete logo as a TikZ vector graphic
% Usage: \isaraLogo[scale=1.0, opacity=1.0, color=black]
\newcommand{\isaraLogo}[1][]{%
  \begin{tikzpicture}[#1, scale=1, baseline=0]
    \drawIsaraLogoShape{black}
  \end{tikzpicture}%
}

% =============================================================================




% -----------------------------------------------------------------------------
% Beamer theme configuration (2025 Modern Style)
% -----------------------------------------------------------------------------
\setbeamertemplate{navigation symbols}{}
\setbeamertemplate{blocks}[rounded][shadow=false]

% Block colors - using new palette
\setbeamercolor{block title}{bg=mvaPrimary,fg=white}
\setbeamercolor{block body}{bg=mvaBackground}

% Alert and example blocks
\setbeamercolor{block title alerted}{bg=mvaAlert,fg=white}
\setbeamercolor{block body alerted}{bg=mvaAlert!5}
\setbeamercolor{block title example}{bg=mvaSuccess,fg=white}
\setbeamercolor{block body example}{bg=mvaSuccess!5}

% Frame elements
\setbeamercolor{frametitle}{fg=mvaPrimary}
\setbeamercolor{framesubtitle}{fg=mvaMuted}
\setbeamercolor{title}{fg=white}
\setbeamercolor{subtitle}{fg=white!90}

% Item colors
\setbeamercolor{itemize item}{fg=mvaAccent}
\setbeamercolor{itemize subitem}{fg=mvaPrimary}
\setbeamercolor{enumerate item}{fg=mvaAccent}
\setbeamercolor{enumerate subitem}{fg=mvaPrimary}

% Structure colors
\setbeamercolor{structure}{fg=mvaPrimary}
\setbeamercolor{alerted text}{fg=mvaAlert}

% -----------------------------------------------------------------------------
% Spacing and Layout
% -----------------------------------------------------------------------------
\setbeamersize{text margin left=1.2cm, text margin right=1.2cm}
\setlength{\parskip}{0.4em}

% Block spacing
\addtobeamertemplate{block begin}{\vskip 0.4em}{}
\addtobeamertemplate{block end}{}{\vskip 0.4em}

% Frame title spacing and styling
\addtobeamertemplate{frametitle}{}{\vspace{0.3em}}
\setbeamerfont{frametitle}{series=\bfseries}

% -----------------------------------------------------------------------------
% Item styling
% -----------------------------------------------------------------------------
\setbeamertemplate{itemize item}{\textbullet}
\setbeamertemplate{itemize subitem}{\textendash}
\setbeamertemplate{enumerate items}[default]

% Font settings (fallback for pdflatex)
\renewcommand{\sfdefault}{lmss}
\sffamily

% -----------------------------------------------------------------------------
% Custom commands
% -----------------------------------------------------------------------------
\newcommand{\academicyear}{}
\newcommand{\otoprule}{\midrule[\heavyrulewidth]}

% Colored highlight boxes
\newcommand{\yellbox}[1]{%
  \tikz[baseline]{\node[inner sep=2pt,fill=mvaAccent!30,text=black,rounded corners=0.5mm,anchor=base]{#1};}}

\newcommand{\questionbox}[1]{%
  \tikz[baseline]{\node[inner sep=2pt,fill=mvaHighlight,text=black,rounded corners=0.5mm,anchor=base]{#1};}}

\newcommand{\rightyellbox}[1]{\begin{flushright}\yellbox{#1}\end{flushright}}
\newcommand{\rightquestionbox}[1]{\begin{flushright}\questionbox{#1}\end{flushright}}

% Course info
\newcommand{\courseurl}{https://misovalko.github.io/mva-ml-graphs.html}
\newcommand{\instructoremail}{michal.valko@inria.fr}

% =============================================================================


% === MERGED SLIDE MACROS ===

% === Macro: MakeNormalSlide ===

\newcommand\MakeNormalSlide{
% =============================================================================
% Normal Slide Footer Template (2025 Modern Edition with Isara Logo)
% =============================================================================

% Clean white background with huge transparent Isara logo watermark
% The logo is always there, subtly, like the foundation from which all content emerges
% Using beautiful TikZ vector graphics - a piece of art
\setbeamertemplate{background canvas}{%
  \begin{tikzpicture}[remember picture,overlay]
    % Pure white background - no gradient to ensure text readability
    \fill[white] (current page.south west) rectangle (current page.north east);

    % Huge Isara logo - centered, visible watermark
    % This represents the foundation, always present and visible
    \node[opacity=0.25,anchor=center,inner sep=0] at (current page.center) {
      \begin{tikzpicture}[scale=25, transform shape]
        \drawIsaraLogoShape{mvaPrimary!50}
      \end{tikzpicture}
    };
  \end{tikzpicture}%
}
\setbeamercolor{toto}{fg=blanc,bg=mvaPrimary}

% Modern minimal footer
\setbeamertemplate{footline}{%
  \begin{beamercolorbox}[wd=\paperwidth,ht=14pt]{toto}
    \hspace{4mm}%
    \raisebox{2.2ex}{\fontsize{8}{10}\selectfont Michal Valko -- Graphs in Machine Learning}%
    \hfill%
    \raisebox{2.2ex}{\fontsize{8}{10}\selectfont\insertframenumber/\inserttotalframenumber\hspace{4mm}\null}
  \end{beamercolorbox}%
}

% Frame title styling
\setbeamertemplate{frametitle}{%
  \vspace{0.3em}%
  \textbf{\insertframetitle}%
  \ifx\insertframesubtitle\empty\else
    \\[0.2em]{\small\textcolor{mvaMuted}{\insertframesubtitle}}%
  \fi
}

}

% === Macro: MakeTitleSlide ===

\newcommand\MakeTitleSlide{
% =============================================================================
% Main Title Slide Template (2025 Modern Edition with Isara Logo)
% =============================================================================
% Used for Lecture 0 (Overview)

% Modern gradient background with Isara logo - "born from it" design
% The logo emerges from the depths, creating the foundation from which content arises
% Using beautiful TikZ vector graphics - a piece of art with individual logo shapes
\setbeamertemplate{background canvas}{%
  \begin{tikzpicture}[remember picture,overlay]
    % Gradient background - deeper, richer
    \shade[top color=mvaPrimary!95,middle color=mvaPrimary!85,bottom color=mvaPrimary!70!black] 
      (current page.south west) rectangle (current page.north east);
    
    % Full Isara logo - large, emerging from bottom-center, creating the foundation
    \node[opacity=0.35,anchor=south,inner sep=0,yshift=-0.5cm] at (current page.south) {
      \begin{tikzpicture}[scale=14, transform shape]
        \drawIsaraLogoShape{white}
      \end{tikzpicture}
    };
    
    % Individual logo shapes scattered artistically across the slide
    % Circles emerging from various positions
    \node[opacity=0.25,anchor=south west,inner sep=0,xshift=1.5cm,yshift=2cm,rotate=45] at (current page.south west) {
      \begin{tikzpicture}[scale=4, transform shape]
        \drawIsaraCircle{white}
      \end{tikzpicture}
    };
    \node[opacity=0.20,anchor=south east,inner sep=0,xshift=-1.5cm,yshift=2.5cm,rotate=-30] at (current page.south east) {
      \begin{tikzpicture}[scale=3.5, transform shape]
        \drawIsaraCircle{white}
      \end{tikzpicture}
    };
    \node[opacity=0.18,anchor=north west,inner sep=0,xshift=2cm,yshift=-1.5cm,rotate=60] at (current page.north west) {
      \begin{tikzpicture}[scale=3, transform shape]
        \drawIsaraCircle{white}
      \end{tikzpicture}
    };
    
    % Main curves floating around
    \node[opacity=0.22,anchor=south west,inner sep=0,xshift=0.5cm,yshift=0.5cm,rotate=-15] at (current page.south west) {
      \begin{tikzpicture}[scale=8, transform shape]
        \drawIsaraMainCurve{white}
      \end{tikzpicture}
    };
    \node[opacity=0.16,anchor=south east,inner sep=0,xshift=-0.5cm,yshift=0.5cm,rotate=12] at (current page.south east) {
      \begin{tikzpicture}[scale=6, transform shape]
        \drawIsaraMainCurve{white}
      \end{tikzpicture}
    };
    \node[opacity=0.18,anchor=north west,inner sep=0,xshift=0.8cm,yshift=-0.8cm,rotate=28] at (current page.north west) {
      \begin{tikzpicture}[scale=7, transform shape]
        \drawIsaraMainCurve{white}
      \end{tikzpicture}
    };
    
    % Swirls creating movement and flow
    \node[opacity=0.20,anchor=west,inner sep=0,xshift=1cm,rotate=25] at (current page.west) {
      \begin{tikzpicture}[scale=5, transform shape]
        \drawIsaraTopSwirl{white}
      \end{tikzpicture}
    };
    \node[opacity=0.18,anchor=east,inner sep=0,xshift=-1cm,rotate=-35] at (current page.east) {
      \begin{tikzpicture}[scale=4.5, transform shape]
        \drawIsaraBottomSwirl{white}
      \end{tikzpicture}
    };
    \node[opacity=0.15,anchor=north,inner sep=0,yshift=-2cm,rotate=45] at (current page.north) {
      \begin{tikzpicture}[scale=4, transform shape]
        \drawIsaraTopSwirl{white}
      \end{tikzpicture}
    };
    
    % Curves adding elegance
    \node[opacity=0.16,anchor=north east,inner sep=0,xshift=-0.8cm,yshift=-0.8cm,rotate=-22] at (current page.north east) {
      \begin{tikzpicture}[scale=5, transform shape]
        \drawIsaraLeftCurve{white}
      \end{tikzpicture}
    };
    \node[opacity=0.14,anchor=south,inner sep=0,yshift=1.5cm,xshift=-3cm,rotate=50] at (current page.south) {
      \begin{tikzpicture}[scale=3.5, transform shape]
        \drawIsaraRightCurve{white}
      \end{tikzpicture}
    };
    
    % Inner swirls for subtle detail
    \node[opacity=0.12,anchor=center,inner sep=0,xshift=-4cm,yshift=1cm,rotate=-40] at (current page.center) {
      \begin{tikzpicture}[scale=3, transform shape]
        \drawIsaraInnerSwirl{white}
      \end{tikzpicture}
    };
    \node[opacity=0.10,anchor=center,inner sep=0,xshift=4cm,yshift=-1cm,rotate=55] at (current page.center) {
      \begin{tikzpicture}[scale=2.5, transform shape]
        \drawIsaraInnerSwirl{white}
      \end{tikzpicture}
    };
    
    % Radial gradient overlay to blend logo into background - creating depth
    \shade[inner color=transparent,outer color=mvaPrimary!75!black] 
      (current page.center) circle (9.5cm);
    % Additional subtle radial from bottom to create "emerging" effect
    \shade[inner color=mvaPrimary!40,outer color=transparent] 
      ([yshift=-2cm]current page.south) ellipse (10cm and 6cm);
  \end{tikzpicture}%
}
\setbeamertemplate{footline}{%
  \hspace{1.8em}\textcolor{white!70}{\null\hfill\textbf{MVA}}\hspace{2em}\null\vspace*{3pt}}

\begin{frame}
% Prominent Isara logo on the first page - TikZ vector graphic
\begin{tikzpicture}[remember picture,overlay]
  % Logo in top-right corner, prominent but elegant
  \node[anchor=north east,inner sep=0,xshift=-3cm,yshift=-1cm] at (current page.north east) {
    \begin{tikzpicture}[scale=2.5, transform shape]
      \drawIsaraLogoShape{white!85}
    \end{tikzpicture}
  };
\end{tikzpicture}

\begin{textblock*}{12cm}(4mm,40mm)
  \textcolor{white}{%
    {\fontsize{22}{26}\selectfont\textbf{Graphs in Machine Learning}}\\[4mm]
    {\fontsize{11}{13}\selectfont Michal Valko}\\[1mm]
    {\fontsize{9}{11}\selectfont\textit{Inria \& ENS Paris-Saclay, MVA}}\\[2mm]
    \ifdefined\misoclasstinytitleslide
      {\fontsize{7}{9}\selectfont\misoclasstinytitleslide}
    \fi
  }
\end{textblock*}
\begin{textblock*}{12cm}(60mm,40mm)
  \includegraphics[width=6.5cm]{../img/peer_communities2.png}
\end{textblock*}
\vspace*{-4pt}
\end{frame}

\MakeNormalSlide

}

% === Macro: MakeTopicSlide ===

\newcommand\MakeTopicSlide{
% =============================================================================
% Topic-Specific Title Slide Template (2025 Modern Edition with Isara Logo)
% =============================================================================
% Usage:
%   \newcommand{\topictitle}{Spectral Clustering}
%   \newcommand{\topicsubtitle}{Theory and Applications}  % optional
%   \newcommand{\topiccredits}{Partially based on material by: ...}  % optional
%   \input{../common/mvamlgraphstopic1slide}

% Modern gradient background with Isara logo - "born from it" design
% The logo emerges from the depths, creating the foundation from which content arises
% Using beautiful TikZ vector graphics - a piece of art with individual logo shapes
\setbeamertemplate{background canvas}{%
  \begin{tikzpicture}[remember picture,overlay]
    % Gradient background - deeper, richer
    \shade[top color=mvaPrimary!95,middle color=mvaPrimary!85,bottom color=mvaPrimary!70!black]
      (current page.south west) rectangle (current page.north east);

    % Full Isara logo - large, emerging from bottom-center, creating the foundation
    \node[opacity=0.35,anchor=south,inner sep=0,yshift=-0.5cm] at (current page.south) {
      \begin{tikzpicture}[scale=14, transform shape]
        \drawIsaraLogoShape{white}
      \end{tikzpicture}
    };

    % Individual logo shapes scattered artistically across the slide
    % Circles emerging from various positions - artistic modern design
    \node[opacity=0.18,anchor=south west,inner sep=0,xshift=1.5cm,yshift=2cm,rotate=45] at (current page.south west) {
      \begin{tikzpicture}[scale=4, transform shape]
        \drawIsaraCircle{white}
      \end{tikzpicture}
    };
    \node[opacity=0.15,anchor=south east,inner sep=0,xshift=-1.5cm,yshift=2.5cm,rotate=-30] at (current page.south east) {
      \begin{tikzpicture}[scale=3.5, transform shape]
        \drawIsaraCircle{white}
      \end{tikzpicture}
    };
    \node[opacity=0.12,anchor=north west,inner sep=0,xshift=2cm,yshift=-1.5cm,rotate=60] at (current page.north west) {
      \begin{tikzpicture}[scale=3, transform shape]
        \drawIsaraCircle{white}
      \end{tikzpicture}
    };

    % Main curves floating around - artistic modern design
    \node[opacity=0.16,anchor=south west,inner sep=0,xshift=0.5cm,yshift=0.5cm,rotate=-15] at (current page.south west) {
      \begin{tikzpicture}[scale=8, transform shape]
        \drawIsaraMainCurve{white}
      \end{tikzpicture}
    };
    \node[opacity=0.12,anchor=south east,inner sep=0,xshift=-0.5cm,yshift=0.5cm,rotate=12] at (current page.south east) {
      \begin{tikzpicture}[scale=6, transform shape]
        \drawIsaraMainCurve{white}
      \end{tikzpicture}
    };
    \node[opacity=0.14,anchor=north west,inner sep=0,xshift=0.8cm,yshift=-0.8cm,rotate=28] at (current page.north west) {
      \begin{tikzpicture}[scale=7, transform shape]
        \drawIsaraMainCurve{white}
      \end{tikzpicture}
    };

    % Swirls creating movement and flow - artistic modern design
    \node[opacity=0.14,anchor=west,inner sep=0,xshift=1cm,rotate=25] at (current page.west) {
      \begin{tikzpicture}[scale=5, transform shape]
        \drawIsaraTopSwirl{white}
      \end{tikzpicture}
    };
    \node[opacity=0.12,anchor=east,inner sep=0,xshift=-1cm,rotate=-35] at (current page.east) {
      \begin{tikzpicture}[scale=4.5, transform shape]
        \drawIsaraBottomSwirl{white}
      \end{tikzpicture}
    };
    \node[opacity=0.10,anchor=north,inner sep=0,yshift=-2cm,rotate=45] at (current page.north) {
      \begin{tikzpicture}[scale=4, transform shape]
        \drawIsaraTopSwirl{white}
      \end{tikzpicture}
    };

    % Curves adding elegance - artistic modern design
    \node[opacity=0.12,anchor=north east,inner sep=0,xshift=-0.8cm,yshift=-0.8cm,rotate=-22] at (current page.north east) {
      \begin{tikzpicture}[scale=5, transform shape]
        \drawIsaraLeftCurve{white}
      \end{tikzpicture}
    };
    \node[opacity=0.10,anchor=south,inner sep=0,yshift=1.5cm,xshift=-3cm,rotate=50] at (current page.south) {
      \begin{tikzpicture}[scale=3.5, transform shape]
        \drawIsaraRightCurve{white}
      \end{tikzpicture}
    };

    % Inner swirls for subtle detail - artistic modern design
    \node[opacity=0.09,anchor=center,inner sep=0,xshift=-4cm,yshift=1cm,rotate=-40] at (current page.center) {
      \begin{tikzpicture}[scale=3, transform shape]
        \drawIsaraInnerSwirl{white}
      \end{tikzpicture}
    };
    \node[opacity=0.08,anchor=center,inner sep=0,xshift=4cm,yshift=-1cm,rotate=55] at (current page.center) {
      \begin{tikzpicture}[scale=2.5, transform shape]
        \drawIsaraInnerSwirl{white}
      \end{tikzpicture}
    };

    % Radial gradient overlay to blend logo into background - creating depth
    \shade[inner color=transparent,outer color=mvaPrimary!75!black]
      (current page.center) circle (9.5cm);
    % Additional subtle radial from bottom to create "emerging" effect
    \shade[inner color=mvaPrimary!40,outer color=transparent]
      ([yshift=-2cm]current page.south) ellipse (10cm and 6cm);
  \end{tikzpicture}%
}
\setbeamertemplate{footline}{%
  \hspace{1.8em}\textcolor{white!70}{\null\hfill\textbf{MVA}}\hspace{2em}\null\vspace*{3pt}}

\begin{frame}
% Modern artistic layout with ISARA logo - ultra-minimal white zone
\begin{tikzpicture}[remember picture,overlay]
  % Single smooth radial gradient from blue to white
  % One continuous gradient zone with microscopic white core
  % Ultra-extended outer radius with ultra-gentle opacity for extremely long, soft transitions

  % Generate single continuous gradient - layers 1-100
  % Smooth transition from blue to white with extreme extended reach across the slide
  \foreach \i in {1,...,100} {
    \pgfmathsetmacro{\radius}{11.0 - \i*0.104}
    \pgfmathsetmacro{\whitePercent}{\i*\i/100}
    \pgfmathsetmacro{\opacity}{0.02 + \i*0.01}
    \shade[inner color=white!\whitePercent!mvaPrimary,outer color=white!\whitePercent!mvaPrimary,opacity=\opacity]
      ([xshift=-0.8cm,yshift=-1.3cm]current page.north east) circle (\radius cm);
  }

  % Microscopic pure white core for logo background
  \fill[white]
    ([xshift=-0.8cm,yshift=-1.3cm]current page.north east) circle (0.6cm);

  % ISARA logo in top-right corner - transparent background
  \node[anchor=north east,inner sep=0,xshift=0.2cm,yshift=-0.3cm] at (current page.north east) {
    
\includegraphics[width=1.8cm]{../common/isara-logo.png}
  };
\end{tikzpicture}

% Modern asymmetric text layout - positioned to avoid logo and graph
\begin{textblock*}{11cm}(6mm,25mm)
  \textcolor{white}{%
    {\fontsize{24}{28}\selectfont\textbf{Graphs in Machine Learning}}\\[4mm]
    {\fontsize{18}{22}\selectfont\textbf{\topictitle}}\\[3mm]
    \ifdefined\topicsubtitle
      {\fontsize{12}{14}\selectfont\topicsubtitle}\\[3mm]
    \fi
    {\fontsize{11}{13}\selectfont Michal Valko}\\[1mm]
    {\fontsize{9}{11}\selectfont\textit{Inria \& ENS Paris-Saclay, MVA}}\\[4mm]
    \ifdefined\topiccredits
      {\fontsize{7}{9}\selectfont\topiccredits}
    \else
      {\fontsize{7}{9}\selectfont\misoclasstinytitleslide}
    \fi
  }
\end{textblock*}
% Graph visualization - positioned artistically to the right
\begin{textblock*}{7cm}(60mm,32mm)
  \includegraphics[width=6.5cm]{../img/peer_communities2.png}
\end{textblock*}
\vspace*{-4pt}
\end{frame}

\MakeNormalSlide

}

% === Macro: MakeLastSlide ===

\newcommand\MakeLastSlide{
% =============================================================================
% Last Slide Template (2025 Modern Edition with Isara Logo)
% =============================================================================

% Modern gradient background with Isara logo - "born from it" design (returning to source)
% The journey completes, content returns to its origin, merging back with the source
% Using beautiful TikZ vector graphics - a piece of art
\setbeamertemplate{background canvas}{%
  \begin{tikzpicture}[remember picture,overlay]
    % Gradient background - deeper, richer
    \shade[top color=mvaPrimary!95,middle color=mvaPrimary!85,bottom color=mvaPrimary!70!black]
      (current page.south west) rectangle (current page.north east);

    % Full Isara logo - large, emerging from bottom-center, creating the foundation
    \node[opacity=0.35,anchor=south,inner sep=0,yshift=-0.5cm] at (current page.south) {
      \begin{tikzpicture}[scale=14, transform shape]
        \drawIsaraLogoShape{white}
      \end{tikzpicture}
    };

    % Individual logo shapes scattered artistically across the slide
    % Circles emerging from various positions - artistic modern design
    \node[opacity=0.18,anchor=south west,inner sep=0,xshift=1.5cm,yshift=2cm,rotate=45] at (current page.south west) {
      \begin{tikzpicture}[scale=4, transform shape]
        \drawIsaraCircle{white}
      \end{tikzpicture}
    };
    \node[opacity=0.15,anchor=south east,inner sep=0,xshift=-1.5cm,yshift=2.5cm,rotate=-30] at (current page.south east) {
      \begin{tikzpicture}[scale=3.5, transform shape]
        \drawIsaraCircle{white}
      \end{tikzpicture}
    };
    \node[opacity=0.12,anchor=north west,inner sep=0,xshift=2cm,yshift=-1.5cm,rotate=60] at (current page.north west) {
      \begin{tikzpicture}[scale=3, transform shape]
        \drawIsaraCircle{white}
      \end{tikzpicture}
    };

    % Main curves floating around - artistic modern design
    \node[opacity=0.16,anchor=south west,inner sep=0,xshift=0.5cm,yshift=0.5cm,rotate=-15] at (current page.south west) {
      \begin{tikzpicture}[scale=8, transform shape]
        \drawIsaraMainCurve{white}
      \end{tikzpicture}
    };
    \node[opacity=0.12,anchor=south east,inner sep=0,xshift=-0.5cm,yshift=0.5cm,rotate=12] at (current page.south east) {
      \begin{tikzpicture}[scale=6, transform shape]
        \drawIsaraMainCurve{white}
      \end{tikzpicture}
    };
    \node[opacity=0.14,anchor=north west,inner sep=0,xshift=0.8cm,yshift=-0.8cm,rotate=28] at (current page.north west) {
      \begin{tikzpicture}[scale=7, transform shape]
        \drawIsaraMainCurve{white}
      \end{tikzpicture}
    };

    % Swirls creating movement and flow - artistic modern design
    \node[opacity=0.14,anchor=west,inner sep=0,xshift=1cm,rotate=25] at (current page.west) {
      \begin{tikzpicture}[scale=5, transform shape]
        \drawIsaraTopSwirl{white}
      \end{tikzpicture}
    };
    \node[opacity=0.12,anchor=east,inner sep=0,xshift=-1cm,rotate=-35] at (current page.east) {
      \begin{tikzpicture}[scale=4.5, transform shape]
        \drawIsaraBottomSwirl{white}
      \end{tikzpicture}
    };
    \node[opacity=0.10,anchor=north,inner sep=0,yshift=-2cm,rotate=45] at (current page.north) {
      \begin{tikzpicture}[scale=4, transform shape]
        \drawIsaraTopSwirl{white}
      \end{tikzpicture}
    };

    % Curves adding elegance - artistic modern design
    \node[opacity=0.12,anchor=north east,inner sep=0,xshift=-0.8cm,yshift=-0.8cm,rotate=-22] at (current page.north east) {
      \begin{tikzpicture}[scale=5, transform shape]
        \drawIsaraLeftCurve{white}
      \end{tikzpicture}
    };
    \node[opacity=0.10,anchor=south,inner sep=0,yshift=1.5cm,xshift=-3cm,rotate=50] at (current page.south) {
      \begin{tikzpicture}[scale=3.5, transform shape]
        \drawIsaraRightCurve{white}
      \end{tikzpicture}
    };

    % Inner swirls for subtle detail - artistic modern design
    \node[opacity=0.09,anchor=center,inner sep=0,xshift=-4cm,yshift=1cm,rotate=-40] at (current page.center) {
      \begin{tikzpicture}[scale=3, transform shape]
        \drawIsaraInnerSwirl{white}
      \end{tikzpicture}
    };
    \node[opacity=0.08,anchor=center,inner sep=0,xshift=4cm,yshift=-1cm,rotate=55] at (current page.center) {
      \begin{tikzpicture}[scale=2.5, transform shape]
        \drawIsaraInnerSwirl{white}
      \end{tikzpicture}
    };

    % Radial gradient overlay to blend logo into background - creating depth
    \shade[inner color=transparent,outer color=mvaPrimary!75!black]
      (current page.center) circle (9.5cm);
    % Additional subtle radial from bottom to create "emerging" effect
    \shade[inner color=mvaPrimary!40,outer color=transparent]
      ([yshift=-2cm]current page.south) ellipse (10cm and 6cm);
  \end{tikzpicture}%
}
\setbeamertemplate{footline}{}

\begin{frame}
% Modern artistic layout with ISARA logo at bottom - ultra-minimal white zone
\begin{tikzpicture}[remember picture,overlay]
  % Single smooth radial gradient from blue to white
  % One continuous gradient zone with microscopic white core
  % Ultra-extended outer radius with ultra-gentle opacity for extremely long, soft transitions

  % Generate single continuous gradient - layers 1-100
  % Smooth transition from blue to white with extreme extended reach across the slide
  \foreach \i in {1,...,100} {
    \pgfmathsetmacro{\radius}{11.0 - \i*0.104}
    \pgfmathsetmacro{\whitePercent}{\i*\i/100}
    \pgfmathsetmacro{\opacity}{0.02 + \i*0.01}
    \shade[inner color=white!\whitePercent!mvaPrimary,outer color=white!\whitePercent!mvaPrimary,opacity=\opacity]
      ([xshift=1.5cm,yshift=2.0cm]current page.south west) circle (\radius cm);
  }

  % Microscopic pure white core for logo background
  \fill[white]
    ([xshift=1.5cm,yshift=2.0cm]current page.south west) circle (0.6cm);

  % ISARA logo in bottom-left corner - transparent background
  \node[anchor=south west,inner sep=0,xshift=0.5cm,yshift=1.0cm] at (current page.south west) {
    
\includegraphics[width=1.8cm]{../common/isara-logo.png}
  };
\end{tikzpicture}

\begin{textblock*}{10cm}(30mm,45mm)
  \textcolor{white}{%
    \raggedright
    {\fontsize{24}{28}\selectfont\textbf{Michal Valko}}\\[4mm]
    {\fontsize{12}{14}\selectfont\ttfamily\instructoremail}\\[2mm]
    {\fontsize{11}{13}\selectfont Inria \& ENS Paris-Saclay, MVA}\\[4mm]
    {\fontsize{10}{12}\selectfont\ttfamily\url{\courseurl}}
  }
\end{textblock*}
\vspace*{-4pt}
\end{frame}

% Reset to normal slide template
\setbeamertemplate{background canvas}[vertical shading][top=blanc,middle=blanc,bottom=blanc]
\setbeamercolor{toto}{fg=blanc,bg=mvaPrimary}

\setbeamertemplate{footline}{%
  \begin{beamercolorbox}[wd=\paperwidth,ht=15.5pt]{toto}
    \hspace{3mm}%
    \includegraphics[width=14mm]{../common/logobastrans}%
    \hspace{1.5cm}%
    \raisebox{2.5ex}{Michal Valko -- Inria}%
    \hfill%
    \raisebox{2.5ex}{\insertframenumber/\inserttotalframenumber\hspace{5mm}\null}
  \end{beamercolorbox}%
}

}
